% GENERATED FILE. Edit canonical lessons, then run npm run generate:books.
% canonical-source-hash: fnv1a64:e16f60119f139319
% canonical-lessons: BN-C259-dor, BN-C259-chhar, BN-C259-phorm, BN-C259-abedon, BN-C259-porichoypotro

\chapter{Rate, Discount, Form, Application, Identity card}
\label{ch:bn-rate-discount-form-application-identity-card}
\hlchaptermodality{259}

\begin{chapteropening}
\textbf{By the end of this chapter:} \emph{I can say a rate, a discount, a form, an application and an identity card.}

Complete the last lesson of chapter 259: I can say a rate, a discount, a form, an application and an identity card.
\end{chapteropening}

\section[\texorpdfstring{\bn{দর} (dôr)}{দর (dôr)}]{\bn{দর} (dôr) — a rate, a price}

\label{lesson:BN-C259-dor}



\begin{quote}
Before the new one: say the Bengali for a battery, then the Bengali for a switch.
\end{quote}

\subsection*{You'll want to know: \bn{দর}}
\textbf{\bn{দর}} — \emph{dôr} — \textquotedblleft{}a rate, a price\textquotedblright{}.

\begin{cousinweb}
The origin of \textbf{\bn{দর}} (\emph{dôr}) is not settled: CDIAL has no entry for it in this sense, and a search by the glosses \textquotedblleft{}rate\textquotedblright{} and \textquotedblleft{}price\textquotedblright{} finds only words of other shapes, such as Bengali \emph{dām} (CDIAL 6622).
\end{cousinweb}

\begin{grammarlens}[title={What you've built}]
One more word for this chapter.
\end{grammarlens}

\subsection*{Your turn}
\begin{itemize}
\raggedright
  \item \emph{Say it:} \emph{dôr}
  \item \emph{Say it:} \emph{dôr}, once more
  \item \emph{Say it:} say \emph{dôr}
  \item \emph{Recall:} say the Bengali for a battery, then the Bengali for a switch, then say \emph{dôr} again
\end{itemize}

\begin{tcolorbox}[breakable,colback=teal!4,colframe=teal!35!black,title={Before you move on}]
What does \bn{দর} mean? (\textquotedblleft{}A rate, a price\textquotedblright{}.) Say it once more.
\end{tcolorbox}
\section[\texorpdfstring{\bn{ছাড়} (chhāṛ)}{ছাড় (chhāṛ)}]{\bn{ছাড়} (chhāṛ) — a discount}

\label{lesson:BN-C259-chhar}



\begin{quote}
Before the new one: say the Bengali for a switch, then the Bengali for a rate.
\end{quote}

\subsection*{You'll want to know: \bn{ছাড়}}
\textbf{\bn{ছাড়}} — \emph{chhāṛ} — \textquotedblleft{}a discount\textquotedblright{}.

\begin{cousinweb}
\textbf{\bn{ছাড়}} (\emph{chhāṛ}), \textquotedblleft{}a discount\textquotedblright{}, appears to belong with the verb \emph{chāṛā}, \textquotedblleft{}to let go\textquotedblright{}, which CDIAL lists under Sanskrit \emph{chardati}, \textquotedblleft{}causes to flow over\textquotedblright{} (CDIAL 4998): \textquotedblleft{}a letting go\textquotedblright{}.
\end{cousinweb}

\begin{grammarlens}[title={What you've built}]
One more word for this chapter.
\end{grammarlens}

\subsection*{Your turn}
\begin{itemize}
\raggedright
  \item \emph{Say it:} \emph{chhāṛ}
  \item \emph{Say it:} \emph{chhāṛ}, once more
  \item \emph{Say it:} say \emph{chhāṛ}
  \item \emph{Recall:} say the Bengali for a switch, then the Bengali for a rate, then say \emph{chhāṛ} again
\end{itemize}

\begin{tcolorbox}[breakable,colback=teal!4,colframe=teal!35!black,title={Before you move on}]
What does \bn{ছাড়} mean? (\textquotedblleft{}A discount\textquotedblright{}.) Say it once more.
\end{tcolorbox}
\section[\texorpdfstring{\bn{ফর্ম} (phôrm)}{ফর্ম (phôrm)}]{\bn{ফর্ম} (phôrm) — a form}

\label{lesson:BN-C259-phorm}



\begin{quote}
Before the new one: say the Bengali for a rate, then the Bengali for a discount.
\end{quote}

\subsection*{You'll want to know: \bn{ফর্ম}}
\textbf{\bn{ফর্ম}} — \emph{phôrm} — \textquotedblleft{}a form\textquotedblright{}.

\begin{cousinweb}
\textbf{\bn{ফর্ম}} (\emph{phôrm}) is a loan from English \emph{form}.
\end{cousinweb}

\begin{grammarlens}[title={What you've built}]
One more word for this chapter.
\end{grammarlens}

\subsection*{Your turn}
\begin{itemize}
\raggedright
  \item \emph{Say it:} \emph{phôrm}
  \item \emph{Say it:} \emph{phôrm}, once more
  \item \emph{Say it:} say \emph{phôrm}
  \item \emph{Recall:} say the Bengali for a rate, then the Bengali for a discount, then say \emph{phôrm} again
\end{itemize}

\begin{tcolorbox}[breakable,colback=teal!4,colframe=teal!35!black,title={Before you move on}]
What does \bn{ফর্ম} mean? (\textquotedblleft{}A form\textquotedblright{}.) Say it once more.
\end{tcolorbox}
\section[\texorpdfstring{\bn{আবেদন} (ābedôn)}{আবেদন (ābedôn)}]{\bn{আবেদন} (ābedôn) — an application}

\label{lesson:BN-C259-abedon}



\begin{quote}
Before the new one: say the Bengali for a discount, then the Bengali for a form.
\end{quote}

\subsection*{You'll want to know: \bn{আবেদন}}
\textbf{\bn{আবেদন}} — \emph{ābedôn} — \textquotedblleft{}an application\textquotedblright{}.

\begin{cousinweb}
\textbf{\bn{আবেদন}} (\emph{ābedôn}) is Sanskrit \emph{āvedana}, \textquotedblleft{}announcing, a petition\textquotedblright{}, taken over whole as a learned loan, a tatsama.
\end{cousinweb}

\begin{grammarlens}[title={What you've built}]
One more word for this chapter.
\end{grammarlens}

\subsection*{Your turn}
\begin{itemize}
\raggedright
  \item \emph{Say it:} \emph{ābedôn}
  \item \emph{Say it:} \emph{ābedôn}, once more
  \item \emph{Say it:} say \emph{ābedôn}
  \item \emph{Recall:} say the Bengali for a discount, then the Bengali for a form, then say \emph{ābedôn} again
\end{itemize}

\begin{tcolorbox}[breakable,colback=teal!4,colframe=teal!35!black,title={Before you move on}]
What does \bn{আবেদন} mean? (\textquotedblleft{}An application\textquotedblright{}.) Say it once more.
\end{tcolorbox}
\section[\texorpdfstring{\bn{পরিচয়পত্র} (pôrichôypôtrô)}{পরিচয়পত্র (pôrichôypôtrô)}]{\bn{পরিচয়পত্র} (pôrichôypôtrô) — an identity card}

\label{lesson:BN-C259-porichoypotro}



\begin{quote}
Before the new one: say the Bengali for a form, then the Bengali for an application.
\end{quote}

\subsection*{You'll want to know: \bn{পরিচয়পত্র}}
\textbf{\bn{পরিচয়পত্র}} — \emph{pôrichôypôtrô} — \textquotedblleft{}an identity card\textquotedblright{}.

\begin{cousinweb}
\textbf{\bn{পরিচয়পত্র}} (\emph{pôrichôypôtrô}) is built from Sanskrit \emph{paricaya}, \textquotedblleft{}acquaintance\textquotedblright{}, and \emph{pattra}, \textquotedblleft{}leaf, letter\textquotedblright{}, both taken over whole: \textquotedblleft{}acquaintance letter\textquotedblright{}.
\end{cousinweb}

\begin{grammarlens}[title={What you've built}]
One more word for this chapter.
\end{grammarlens}

\subsection*{Your turn}
\begin{itemize}
\raggedright
  \item \emph{Say it:} \emph{pôrichôypôtrô}
  \item \emph{Say it:} \emph{pôrichôypôtrô}, once more
  \item \emph{Say it:} say \emph{pôrichôypôtrô}
  \item \emph{Recall:} say the Bengali for a form, then the Bengali for an application, then say \emph{pôrichôypôtrô} again
\end{itemize}

\begin{tcolorbox}[breakable,colback=teal!4,colframe=teal!35!black,title={Before you move on}]
What does \bn{পরিচয়পত্র} mean? (\textquotedblleft{}An identity card\textquotedblright{}.) Say it once more.
\end{tcolorbox}
